\documentclass[11pt]{article}
\usepackage[margin=1in]{geometry}
\usepackage{amsmath,amssymb}
\usepackage{hyperref}
\usepackage{times}

\begin{document}
\thispagestyle{empty}

\noindent
\textbf{Kavya Aggarwal} \\
Department of Computer Science and Engineering \\
GITAM (Deemed to be University) \\
Hyderabad, Telangana 502329, India \\
Email: \texttt{kaggarwa@gitam.in} \\[1.5em]

\noindent
August 28, 2026 \\[1.5em]

\noindent
\textbf{To:} \\
The Editor-in-Chief \\
\textit{IEEE Transactions on Learning Technologies} \\[1.5em]

\noindent
\textbf{Subject: Submission of Original Research Manuscript -- ``Personalizing Open-Domain Video Curricula: A Hybrid Closed-Loop Framework Combining Contextual Bandits, Difficulty-Conditioned BKT, and Interaction Telemetry''} \\[1.2em]

\noindent
Dear Editor-in-Chief and Editorial Board Members,

\vspace{0.8em}
\noindent
I am pleased to submit our original research manuscript titled \textbf{``Personalizing Open-Domain Video Curricula: A Hybrid Closed-Loop Framework Combining Contextual Bandits, Difficulty-Conditioned BKT, and Interaction Telemetry''} for consideration for publication in \textit{IEEE Transactions on Learning Technologies}.

\vspace{0.8em}
\noindent
\textbf{Background and Contribution:} \\
Video-based learning in Massive Open Online Courses (MOOCs) remains predominantly passive, leading to rapid cognitive decay and an absence of individualized learning paths. While traditional Intelligent Tutoring Systems (ITS) provide structured adaptation, they suffer from severe content-authoring bottlenecks that prevent scaling across open-domain video libraries. In this manuscript, we present \textbf{CogniLoop}, an integrated, closed-loop adaptive learning system that combines:
\begin{enumerate}
    \item A transcript-grounded Retrieval-Augmented Generation (RAG) vector pipeline coupled with an Assessment Validation Layer to eliminate hallucination during assessment generation;
    \item A Difficulty-Conditioned Bayesian Knowledge Tracing (DC-BKT) engine that modulates guess and slip parameters across item difficulty tiers;
    \item A Thompson Sampling Contextual Bandit policy (Beta-distributed posterior sampling with continuous fractional updates) for optimal difficulty routing;
    \item An unsupervised $K$-Means telemetry classifier ($K=3$) that extracts video interaction micro-patterns.
\end{enumerate}

\vspace{0.8em}
\noindent
\textbf{Validation and Statistical Rigor:} \\
Under a controlled algorithmic simulation protocol across $N=60$ student profiles spanning Fast, Standard, and Slow learner archetypes, our framework achieves a statistically significant improvement in Simulated Normalized Learning Gain (NLG) over static sequential video viewing ($\Delta \text{NLG} = +0.541$, Welch's $t = 3.768, p = 4.73 \times 10^{-4} < 0.001$, Cohen's $d = 0.973$, Mann-Whitney $U = 709.5$). A multi-parameter sensitivity analysis ($\pm 0.05$ parameter perturbations across 100 runs) and a six-configuration component ablation study quantitatively confirm that contextual difficulty routing ($\Delta \text{NLG} = -0.294$ when fixed) and semantic retrieval grounding ($28.4\%$ context mismatch without RAG) are critical drivers of system efficacy.

\vspace{0.8em}
\noindent
\textbf{Originality and Conflict of Interest Declarations:}
\begin{itemize}
    \item This manuscript represents original research and has not been published previously, nor is it currently under consideration for publication elsewhere.
    \item The author has no competing financial or non-financial conflicts of interest regarding this work.
    \item The complete open-source codebase, reproducible simulation scripts, and Docker containerization files are maintained at: \url{https://github.com/OmegaKavya/CogniLoop}.
\end{itemize}

\vspace{1.0em}
\noindent
Thank you very much for your time and consideration of our work. I look forward to receiving the reviewers' constructive feedback.

\vspace{1.8em}
\noindent
Sincerely, \\[1.5em]
\textbf{Kavya Aggarwal} \\
Department of Computer Science and Engineering \\
GITAM (Deemed to be University), Hyderabad, India \\
Email: \texttt{kaggarwa@gitam.in}

\end{document}
